\section{Limitations}
\label{sec:limitations}

\subsection{Exploratory model selection}
We compare heads and pooling schemes on the same five grouped folds that we use
for reporting. Selecting the best mean-fold configuration makes its reported
cross-validation performance optimistic due to hyperparameter and
model-selection leak. Group-bootstrap intervals condition on the selected
predictions and do not account for that selection process. A future confirmatory
evaluation requires nested grouped cross-validation or a fully locked external
lineage holdout.

\subsection{Cohort and target construction}
The dataset is extremely small (300 sequences total across 113 lineage groups),
making downstream models highly susceptible to overfitting and limiting the
statistical power of our architectural comparisons. Only variable heavy chains
($V_H$) enter the models, so the analysis completely neglects light-chain ($V_L$) contributions and
heavy-light pairing, which are biochemically known to be critical for the
folding, stability, and neutralization breadth of many HIV-1 bnAbs. Furthermore,
CATNAP narrow/weak antibodies are not equivalent to arbitrary non-bnAbs, and the
analysis treats neutralization breadth as a simple, unweighted fraction over viral
panels that vary in size and strain composition, without accounting for
laboratory-specific or panel-level effects.

\subsection{Grouping assumptions}
While metadata-derived clonal lineage is the strongest available grouping key,
patient-level fallback groups may merge unrelated antibodies from the same
donor, while singletons may mask unrecorded genetic relationships. The
sequence-identity clustering heuristic ($\geq 80\%$ identity) supports
separation in aggregate but cannot biologically prove complete developmental
independence or rule out convergent evolution.

\subsection{Adaptation and ESM3 constraints}
We evaluate the LoRA domain adaptation experiment on a single dataset with one
adapter configuration and one training run. Because the adaptation is
unsupervised (using unlabeled OAS sequences), the model may align
representations toward general repertoire characteristics at the expense of
capturing specific somatic hypermutation features required for HIV-1 gp120/gp41
binding. Additionally, using a 1.4B parameter model like ESM3 introduces
substantial computational overhead for inference and adaptation compared to
simpler sequence-based models, and its non-commercial licensing constraints
limit downstream translation.

\subsection{Interpolation and VAE constraints}
The per-residue VAE treats residues as independent training examples and lacks
an antibody-level encoder. By treating residues independently, the model ignores
sequence order, tertiary structural contacts, and global cooperative folding
constraints. This constitutes a severe biophysical simplification. Furthermore,
the VAE failed our round-trip sequence decoding test-gate, meaning the latent
spaces learned at the residue level cannot be reliably mapped back to
constructible antibody sequences. Any claims of sequence design or maturation
interpolation are therefore speculative and physically unvalidated.

\subsection{Reproducibility boundary}
Our verification script validates pre-computed predictions and fails on cohort
or alignment mismatches, but the massive embedding caches and LoRA adapter
weights come from external GPU clusters (Google Colab), and the local
verification run does not recreate them. The saved prediction and evaluation
artifacts therefore reproduce the paper's claims, but a full, local,
from-scratch retraining pipeline does not.
